\section{Evaluation Results}
\label{sec:evaluation-results}

This section is organised according to the three research questions defined in
Section~\ref{subsec:evaluation-rqs}. Quantitative results are reported for all three
research questions. The RQ1 and RQ2 results are drawn from the persisted run, gate,
approval, and rejection records produced by repeated executions of the hybrid pipeline
orchestrator (\texttt{scripts/run\_hybrid\_pipeline.py}) between 28~June~2026 and
9~July~2026, comprising 17 recorded pipeline invocations. This evaluation corpus was
accumulated during iterative development rather than a single controlled experimental
campaign, so the RQ1 and RQ2 counts below are reported as an initial operational sample
rather than a statistically powered trial.

\subsection{RQ1: Operational Feasibility and Reliability}
\label{subsec:results-rq1}

This experiment concerns execution of the complete SysSecOps pipeline for the pass,
review, reject, and degraded-component scenarios described in
Section~\ref{subsec:evaluation-rq1}. Its reporting structure distinguishes a valid
security rejection from an unexpected execution failure. A rejected change is classified
as a completed run when the pipeline produces a valid gate report and correctly prevents
deployment. The degraded-component scenario additionally distinguishes graceful
degradation, in which an unavailable assessment component or credential is recorded as an
explicit status, from an unhandled termination of the workflow.

The evidence base for this experiment is the set of hybrid-pipeline invocations recorded
in the run logs, cross-referenced with the corresponding gate, approval, and rejection
records. Of 17 recorded invocations, 16 produced a persisted run summary; the remaining
invocation wrote only a log header before terminating without a summary and is counted
as a run that did not reach a valid terminal decision. Each completed invocation
independently gates a CDK
branch, an Ansible branch, or both, yielding 24 branch-level gate decisions (9 CDK, 15
Ansible) across the 16 completed runs. Table~\ref{tab:rq1-results-frame} reports the
observed counts; the completion rate is calculated using
Equation~\ref{eq:pipeline-completion-rate}.

\begin{table}[ht]
\centering
\caption{Observed results for the end-to-end pipeline evaluation.}
\label{tab:rq1-results-frame}
\begin{tabular}{|L{5.9cm}|C{3.0cm}|L{4.7cm}|}
\hline
\textbf{Measure} & \textbf{Result} & \textbf{Evidence source} \\
\hline
Total pipeline runs & $17$ & Execution records \\
\hline
Runs reaching a valid terminal decision & $16$ & Gate reports \\
\hline
End-to-end completion rate & $94.1\%$ ($16/17$) & Calculated from run records \\
\hline
Pass, review, and reject outcomes (branch-level, $n=24$) & $18$ / $4$ / $2$ & Persisted decision records \\
\hline
Unexpected pipeline failures & $2$ & Stage results and execution logs \\
\hline
Rejected changes incorrectly deployed & $0$ of $2$ & Gate and deployment records \\
\hline
Review cases executed without approval & $0$ of $4$ & Approval and deployment records \\
\hline
Degraded-component runs reaching a terminal state & $5$ of $5$ ($100\%$) & Scanner status and gate reports \\
\hline
\end{tabular}
\end{table}

The table shows that 16 of the 17 recorded invocations reached the expected terminal
state (pass, review, or reject), for an end-to-end completion rate of $94.1\%$. The one
run that did not complete terminated after writing only its log header, before any
branch produced a gate report; the orchestrator's exception path was not captured, so no
further cause is retained. A distinct execution-stage failure was also observed within a
completed run: one Ansible branch passed its security gate (score $0$) but the
subsequent deployment step returned a non-zero exit code, giving the terminal status
\texttt{deploy\_failed}. Both this case and the incomplete invocation are counted as
unexpected pipeline failures, since neither resulted from a security-gate rejection. Both
of the two reject decisions were withheld from deployment, and each has a matching
rejection record; no rejected change was deployed. Of the four review decisions, three
remained unapproved and undeployed, and one was deployed after an explicit approval
record was written by the reviewer, so no review case was deployed without a matching
approval. All five CDK branches for which the Infracost cost-analysis
scanner reported an error (a timeout or an unexpected response schema) still produced a
valid pass, review, or reject decision, indicating that the gate degrades gracefully when
this optional component is unavailable rather than terminating the run. Because this
sample was accumulated during iterative development rather than a single controlled
campaign, and because the rejection and degraded-component cases are few in number, these
counts should be read as an initial operational indication rather than a statistically
robust estimate of pipeline reliability.

\subsection{RQ2: Hybrid Infrastructure Applicability}
\label{subsec:results-rq2}

\begin{table}[htb!]
\centering
\caption{Observed results for the existing virtual-machine evaluation.}
\label{tab:rq2-results-frame}
\begin{tabular}{|L{5cm}|C{5cm}|L{4cm}|}
\hline
\textbf{Evaluation checkpoint} & \textbf{Result} & \textbf{Verification evidence} \\
\hline
Network reachability & Confirmed for $6/7$ deploy attempts; undetermined for $1$ & Connectivity test \\
\hline
Authenticated Ansible access & Confirmed for $6/7$ deploy attempts; undetermined for $1$ & Ansible connection result \\
\hline
Playbook syntax validation & Passed in $15/15$ attempts & Validation output \\
\hline
Security assessment and gate decision & Pass in $15/15$ attempts (score $0$) & Scanner output and gate report \\
\hline
Enforcement of the gate before execution & Gate executed prior to deploy in $15/15$ attempts & Decision and execution records \\
\hline
Application of the requested configuration & Succeeded in $6/7$ deploy attempts & Ansible execution result \\
\hline
Verification of the target state & Not captured by retained artifacts & Direct observation of the virtual machine \\
\hline
Repeated-run state consistency & Not directly evidenced by retained artifacts & Subsequent Ansible execution result \\
\hline
Operational and audit record creation & Created for $16/16$ completed runs & Status and audit records \\
\hline
\end{tabular}
\end{table}

This experiment concerns integration of an existing Linux virtual machine through the
Ansible path of the prototype, reached over a mesh-VPN connection and treated as the
private side of the hybrid environment. The same risk engine and decision thresholds
applied to public-cloud changes govern this private-side target. Evidence is collected at
the connectivity, assessment, decision, execution, state-verification, and
operational-recording stages. The on-premises target is addressed in the checked-in
inventory (\texttt{examples/hybrid-webapp-demo/ansible/inventory.ini}) by its Tailscale
mesh address, consistent with the architecture in Section~\ref{subsec:arch-components}.
Evidence is drawn from the 15 Ansible-branch invocations within the hybrid-pipeline run
history reported in Section~\ref{subsec:results-rq1}, of which 7 proceeded to an actual
deployment attempt against the target host. Table~\ref{tab:rq2-results-frame} reports the
observed outcomes; two checkpoints for which the retained artifacts contain no direct
evidence are marked accordingly rather than assigned an unsupported result.



The checkpoint records show that every recorded Ansible-branch attempt produced a
security assessment with zero findings and a pass decision, and that six of the seven
attempts that proceeded to a deployment step completed successfully against the
mesh-VPN-connected host, while one returned a non-zero exit code whose specific cause
(connectivity, authentication, or a task-level condition) cannot be distinguished from
the retained artifacts. Because every observed Ansible-branch decision in this sample was
a pass, the experiment demonstrates that the gate executes ahead of deployment but does
not yet provide direct evidence that a review or reject decision on this path is enforced
in the same way as on the CDK path; that enforcement is inferred from the shared decision
function (Section~\ref{subsec:rs-additive}) rather than observed directly for the
on-premises target. Post-deployment state verification and repeated-run idempotency were
not captured by the retained artifacts and are reported as evidence gaps rather than as
results, motivating a dedicated verification step in a future execution of this
experiment.

\subsection{RQ3: Unified Risk-Evaluation Capability}
\label{subsec:results-rq3}

The available RQ3 evidence comprises the held-out performance of the
logistic-regression model and an examination of its fitted coefficients. This analysis
assesses whether features obtained from multiple static analysers provide a useful
learned signal for the unified risk-scoring mechanism. The reported results concern
detector-visible insecure code and do not represent the detection of every semantic
vulnerability.

\subsubsection{Dataset Composition}
\label{subsubsec:results-dataset}

The corpus was constructed from two classes of Python files. Approximately 80 positive
examples were selected from SecurityEval. Vulnerability categories that could not be
represented by the configured Bandit and Semgrep features were excluded from the model
corpus. The negative reference class contained approximately 80 files sampled from
maintained Python projects, including \texttt{click}, \texttt{rich}, and
\texttt{typer}. To reduce the difference in scale between the short SecurityEval
examples and project source files, negative examples were restricted to files containing
more than 20 and fewer than 300 non-blank, non-comment lines.

The negative examples are referred to as \emph{reference secure} files rather than
formally verified secure files. Maintained project code can still contain unknown
vulnerabilities, and the assigned negative labels therefore represent a practical
reference assumption. Similarly, excluding vulnerability categories outside the
observable scope of the selected analysers narrows the interpretation of the results.
The reported performance applies to the retained, detector-visible subset and is not an
estimate of performance over the complete SecurityEval dataset.

The held-out test partition contained 31 files: 16 negative samples and 15 positive
samples. The near-balanced distribution limited the influence of class frequency on the
reported accuracy.

\subsubsection{Classification Performance}
\label{subsubsec:results-classification}

Table~\ref{tab:rq3-classification-results} presents the aggregate performance of the
logistic-regression classifier on the held-out test partition. The classifier achieved
an accuracy of $0.8710$ and a ROC--AUC of $0.8792$. Precision, recall, and F1-score for
the positive class were all $0.8667$.

\begin{table}[ht]
\centering
\caption{Logistic-regression performance on the held-out test partition.}
\label{tab:rq3-classification-results}
\begin{tabular}{|L{5.5cm}|C{3.5cm}|}
\hline
\textbf{Measure} & \textbf{Result} \\
\hline
Accuracy & $0.8710$ \\
\hline
Precision & $0.8667$ \\
\hline
Recall & $0.8667$ \\
\hline
F1-score & $0.8667$ \\
\hline
ROC--AUC & $0.8792$ \\
\hline
Test samples & $31$ \\
\hline
\end{tabular}
\end{table}

The confusion matrix is presented in Table~\ref{tab:rq3-confusion-matrix}. Of the 16
negative samples, 14 were classified correctly and two were incorrectly classified as
insecure. Of the 15 positive samples, 13 were identified correctly and two were
incorrectly classified as reference secure. The model therefore produced four errors
across the 31 held-out samples.

\begin{table}[ht]
\centering
\caption{Confusion matrix for the held-out test partition.}
\label{tab:rq3-confusion-matrix}
\begin{tabular}{|L{4.7cm}|C{3.4cm}|C{3.4cm}|}
\hline
 & \textbf{Predicted negative} & \textbf{Predicted positive} \\
\hline
\textbf{Actual negative} & $14$ & $2$ \\
\hline
\textbf{Actual positive} & $2$ & $13$ \\
\hline
\end{tabular}
\end{table}

Table~\ref{tab:rq3-class-report} gives the class-specific results. The negative class
obtained precision, recall, and F1-score values of approximately $0.88$, while the
positive class obtained approximately $0.87$ for all three measures. The macro and
weighted averages were both approximately $0.87$, indicating that performance was
similar across the two classes in the held-out partition.

\begin{table}[ht]
\centering
\caption{Class-specific classification results.}
\label{tab:rq3-class-report}
\begin{tabular}{|L{3.6cm}|C{2.5cm}|C{2.5cm}|C{2.5cm}|C{2.2cm}|}
\hline
\textbf{Class} & \textbf{Precision} & \textbf{Recall} & \textbf{F1-score} &
\textbf{Support} \\
\hline
Negative ($0$) & $0.88$ & $0.88$ & $0.88$ & $16$ \\
\hline
Positive ($1$) & $0.87$ & $0.87$ & $0.87$ & $15$ \\
\hline
Macro average & $0.87$ & $0.87$ & $0.87$ & $31$ \\
\hline
Weighted average & $0.87$ & $0.87$ & $0.87$ & $31$ \\
\hline
\end{tabular}
\end{table}

The equal positive-class precision and recall resulted from the symmetric number of
false-positive and false-negative predictions. From an operational perspective, the two
false negatives are particularly relevant because they represent labelled insecure files
that would receive a lower learned risk contribution. The two false positives have a
different consequence: they may increase the score of reference negative code and cause
unnecessary review, but they do not directly permit an insecure file to bypass the
learned component.

\subsubsection{Analysis of Model Coefficients}
\label{subsubsec:results-coefficients}

Table~\ref{tab:rq3-coefficients} presents the fitted logistic-regression coefficients in
descending order. Because the features were standardised before training, coefficient
magnitudes provide an indication of their relative influence within this fitted model. A
positive coefficient indicates that an increase in the corresponding feature increases
the estimated log-odds of the insecure class, with the remaining features held constant.

\begin{table}[ht]
\centering
\small
\caption{Coefficients of the trained logistic-regression model.}
\label{tab:rq3-coefficients}
\begin{tabular}{|L{6.4cm}|C{3.5cm}|}
\hline
\textbf{Feature} & \textbf{Coefficient} \\
\hline
\texttt{total\_semgrep} & $1.136296$ \\
\hline
\texttt{bandit\_conf\_medium} & $1.003949$ \\
\hline
\texttt{semgrep\_high} & $0.914679$ \\
\hline
\texttt{semgrep\_medium} & $0.720708$ \\
\hline
\texttt{total\_bandit} & $0.610932$ \\
\hline
\texttt{bandit\_medium} & $0.454592$ \\
\hline
\texttt{bandit\_high} & $0.435620$ \\
\hline
\texttt{bandit\_low} & $0.261740$ \\
\hline
\texttt{bandit\_conf\_high} & $0.138912$ \\
\hline
\texttt{bandit\_conf\_low} & $0.083852$ \\
\hline
\texttt{semgrep\_low} & $0.067123$ \\
\hline
\end{tabular}
\end{table}

The total number of Semgrep findings had the largest positive coefficient
($1.136296$), followed by the number of medium-confidence Bandit findings
($1.003949$) and high-severity Semgrep findings ($0.914679$). This ordering indicates
that the fitted model relied most strongly on the overall Semgrep signal and on
medium-to-high-strength findings from the two analysers. Low-severity and low-confidence
features had smaller positive coefficients and therefore contributed less strongly to
the predicted probability.

These coefficients must be interpreted with caution. The total-finding features are
mathematically related to the corresponding severity counts, which introduces correlated
predictors. The coefficients describe the behaviour of the fitted model and do not
establish that one feature has a causal relationship with software insecurity.

\subsubsection{Per-File Risk Records}
\label{subsubsec:results-per-file}

In addition to the aggregate classification measures, the corpus is reprocessed through
the trained component to obtain a predicted probability and risk contribution for each
file. The resulting machine-readable artifact is archived with the project repository
for inspection and reproducibility. Each record contains the sample identifier, dataset
source, assigned label, extracted feature values, predicted probability, predicted
class, and learned risk contribution. When the complete set of gate inputs is available,
the record also contains the component scores, total score, and resulting deployment
decision.

Table~\ref{tab:rq3-per-file-fields} defines the fields retained in this supplementary
artifact. The aggregate distribution of per-file probabilities and risk scores is
reported after the re-execution records have been finalised.

\begin{table}[ht]
\centering
\small
\caption{Fields retained in the supplementary per-file risk artifact.}
\label{tab:rq3-per-file-fields}
\begin{tabular}{|L{4.2cm}|L{9.2cm}|}
\hline
\textbf{Field} & \textbf{Description} \\
\hline
Sample identifier & Unique identifier or relative path of the evaluated Python file \\
\hline
Source & SecurityEval or the originating reference project \\
\hline
Assigned label & Positive ($1$) or negative reference ($0$) class \\
\hline
Feature vector & Bandit and Semgrep counts supplied to the trained model \\
\hline
Model output & Predicted probability, predicted class, and learned risk contribution \\
\hline
Gate output & Component scores, total score, and decision when all gate inputs are available \\
\hline
\end{tabular}
\end{table}

Table~\ref{tab:rq3-gate-results-frame} presents the observed outcomes of a controlled
gate-decision check against the implemented scoring logic. The check evaluates the two
decision boundaries directly and confirms that semantically equivalent findings reported
by multiple tools are deduplicated into a single retained finding.

\begin{table}[ht]
\centering
\small
\caption{Observed results for the controlled gate-decision evaluation.}
\label{tab:rq3-gate-results-frame}
\begin{tabular}{|L{4.6cm}|C{2.5cm}|C{2.5cm}|C{2.7cm}|}
\hline
\textbf{Case} & \textbf{Expected} & \textbf{Observed} & \textbf{Agreement} \\
\hline
Score below or equal to pass threshold & Pass & Pass at $S=20$ & Yes \\
\hline
Score immediately above pass threshold & Review & Review at $S=21$ & Yes \\
\hline
Score equal to review upper boundary & Review & Review at $S=80$ & Yes \\
\hline
Score above reject threshold & Reject & Reject at $S=81$ & Yes \\
\hline
Duplicate finding reported by multiple tools & Counted once & One deduplicated finding retained & Yes \\
\hline
\end{tabular}
\end{table}

\subsubsection{Summary of Risk-Model Findings}
\label{subsubsec:results-rq3-summary}

Within the detector-visible and filtered evaluation corpus, the logistic-regression
component achieved an accuracy of $0.8710$, an F1-score of $0.8667$, and a
ROC--AUC of $0.8792$. The similar class-specific measures indicate that the result was
not obtained by favouring only one class. The coefficient analysis further shows that
the model incorporated signals from both Bandit and Semgrep, with the strongest fitted
influence associated with Semgrep finding volume, medium-confidence Bandit findings, and
high-severity Semgrep findings.

The observed performance indicates that a learned multi-tool signal can contribute
discriminatory information to the proposed risk-scoring mechanism within the evaluated
corpus. It does not establish complete vulnerability detection, because the corpus
excludes vulnerability categories outside the configured analysers' detection scope and
uses maintained project files as a reference negative class. Evaluation of the complete
deployment-decision mechanism additionally requires the per-file gate records and the
controlled score-boundary cases defined in
Section~\ref{subsubsec:evaluation-gate-decision}.